% ============================================================================
% Chapter 15 --- Sizing reference points
% ----------------------------------------------------------------------------
% Purpose : indicative orders of magnitude per layer for a typical
%           research-intensive institution. Methodological chapter — the figures are ranges,
%           the assumptions are explicit, the figures are not measurements.
% Page target: ~3 pages.
% Dependencies: Ch. 6--14 (each layer's reference instantiation).
% ----------------------------------------------------------------------------
% Decisions registered in _coherence-EN.md:
%   D15.1 reference profile: research-intensive institution, ~25k users, ~1.5k staff,
%         ~500 active research projects
%   D15.2 all figures are indicative ranges, not measurements
%   D15.3 the chapter does NOT replace local sizing exercises
% ============================================================================

\chapter{Sizing reference points}
\label{ch:sizing}

\begin{caveatbox}[Status of this chapter]
\noindent The figures presented in this chapter are \emph{indicative
orders of magnitude} for a reference institution profile defined in
\S\ref{sec:sz-method}. They are not measurements, and they do not
substitute for the institution's own sizing exercises. Their
purpose is to give the reader a sense of the scale at which the
architecture has been considered, and to flag the points at which
the architecture may need reshaping. Any institutional procurement
or capacity-planning decision should be informed by direct
measurement of the institution's actual workload, not by the
ranges below.
\end{caveatbox}

%-----------------------------------------------------------------------------
\section{Method and assumptions}
\label{sec:sz-method}

\paragraph{Reference profile.} The figures in this chapter assume a
research-intensive university with the following indicative
characteristics:

\begin{itemize}[leftmargin=1.5em,nosep]
  \item approximately twenty-five thousand active subjects (staff,
    students, contractors, visiting researchers combined);
  \item approximately one thousand five hundred staff members with
    administrative or research roles requiring elevated access;
  \item on the order of five hundred active research projects with
    distinct workspace requirements;
  \item a mix of institutional, organisational-unit-level, and
    project-level workloads, including a non-trivial research-data
    footprint.
\end{itemize}

\paragraph{Scope of the figures.} For each layer, two kinds of
figures are given: \emph{steady-state} ranges describing the load
the layer is expected to sustain, and \emph{peak} ranges describing
the load the layer should be provisioned to absorb. Peak figures are
not multiples of steady-state; they are the result of distinct
phenomena (start-of-semester authentication surges, conference
deadline storage pressure, end-of-year research-data ingestion).

\paragraph{Sources of variation.} Actual figures at any institution
are influenced by factors the reference profile does not capture:
the proportion of remote versus on-campus subjects, the depth of
the research-computing footprint, the mix of teaching modalities,
the institutional federation activity, the openness posture toward
visiting researchers. The ranges below would typically span the
P25--P75 of institutions in the broad reference profile, with
substantial overflow on either side.

%-----------------------------------------------------------------------------
\section{Identity layer sizing}
\label{sec:sz-identity}

\begin{itemize}[leftmargin=1.5em,nosep]
  \item Authentications per second, steady state: indicatively
    \emph{a few to tens}; peak (start-of-day, post-incident
    re-authentication storms): indicatively \emph{hundreds}.
  \item Concurrent active sessions: indicatively \emph{thousands};
    peak: indicatively \emph{a few thousand}.
  \item Federation transactions per day (\loc{identity-federation-profile}
    inbound and outbound combined): indicatively \emph{tens to hundreds of
    thousands}.
  \item SCIM provisioning operations per day: indicatively
    \emph{thousands during steady state}; peak around academic-
    year boundaries can be substantially higher for a bounded
    period.
\end{itemize}

\noindent The principal sizing variable is not the user count but
the authentication frequency, which is driven by session-lifetime
conventions: institutions with short sessions and frequent
re-challenge have peak rates an order of magnitude higher than
institutions with long sessions.

%-----------------------------------------------------------------------------
\section{Policy layer sizing}
\label{sec:sz-policy}

\begin{itemize}[leftmargin=1.5em,nosep]
  \item Decisions per second, steady state: indicatively
    \emph{tens to low hundreds}; peak: indicatively \emph{a few
    hundred}.
  \item Decision latency budget: indicatively \emph{below ten
    milliseconds at the 95th percentile} for online decisions
    (network attachment, application authorisation); higher
    budgets are admissible for asynchronous decisions (workload
    admission, batch operations).
  \item Bundle refresh cadence: indicatively \emph{minutes}, with
    the ability to push out-of-band for urgent policy changes.
\end{itemize}

\noindent Policy-decision rate is dominated by the integration
choices, not by the user population: a single authorisation check
embedded in a frequently-called API endpoint can produce more
decisions per second than the rest of the institution combined.

%-----------------------------------------------------------------------------
\section{Network layer sizing}
\label{sec:sz-network}

\begin{itemize}[leftmargin=1.5em,nosep]
  \item Campus backbone capacity: indicatively in the range of
    \emph{40--400~Gb/s aggregate}, with research-intensive
    institutions toward the upper end.
  \item Science DMZ throughput, where present: indicatively
    \emph{at the same order as the backbone}, with the explicit
    architectural choice that the Science DMZ is not penalised by
    enterprise filtering.
  \item NAC transactions per day (wired plus wireless attachment
    and re-attachment): indicatively \emph{hundreds of thousands
    to low millions}.
  \item IPFIX flow records per day: indicatively \emph{hundreds of
    millions}, depending on flow-export sampling rate and on
    research traffic volume.
\end{itemize}

%-----------------------------------------------------------------------------
\section{Audit and observability layer sizing}
\label{sec:sz-audit}

\begin{itemize}[leftmargin=1.5em,nosep]
  \item Events per day, steady state: indicatively \emph{tens of
    millions to hundreds of millions}, depending on the
    instrumentation density.
  \item Storage volume per day, steady state: indicatively
    \emph{tens to low hundreds of gigabytes} after standard
    compression; peak can substantially exceed steady-state during
    incidents.
  \item Retention scenario (illustrative): hot-tier (indexed,
    searchable) for ninety days; warm-tier for one year;
    cold-tier (object storage, retrievable) for the period required
    by regulatory and audit obligations, typically several years.
\end{itemize}

\noindent The audit-layer figures are the most variable in this
chapter: an institution that has not yet operationalised the
canonical schema of \trchap{ch:observability}{} will produce a
small fraction of these volumes; an institution that instruments
aggressively without volume-management discipline can produce
multiples of them.

%-----------------------------------------------------------------------------
\section{FTE indicators per layer}
\label{sec:sz-fte}

The table below consolidates the indicative operational envelopes
introduced in the layer chapters. Each row is a range, not a point
estimate; rows are not additive (engineers are typically shared
across layers within the same operational team).

\begin{table}[H]
\centering\small
\caption{Indicative operational envelopes per layer (engineer-equivalents,
shared with adjacent duties). Ranges are not additive.}
\label{tab:sz-fte}
\begin{tabularx}{\textwidth}{%
  >{\hsize=0.75\hsize\linewidth=\hsize\bfseries\raggedright\arraybackslash}X%
  >{\hsize=0.45\hsize\linewidth=\hsize\centering\arraybackslash}X%
  >{\hsize=0.80\hsize\linewidth=\hsize\raggedright\arraybackslash}X%
}
\toprule
Layer & Indicative range & Notes \\
\midrule
Identity (Ch. 6)         & 2--4   & Reference instantiation; subject lifecycle automation reduces the upper bound. \\
Policy (Ch. 7)           & 1--2   & Skill investment substantial relative to the headcount. \\
Network (Ch. 8)          & 3--5   & Largest envelope; reflects the integrated stack. \\
Compute (Ch. 9)          & 4--7   & Sustained engineering investment for OpenStack + Kubernetes. \\
Storage (Ch. 10)         & 2--4   & Ceph operations require dedicated familiarity. \\
Endpoint (Ch. 11)        & 3--6   & Multi-OS coverage drives the count. \\
Observability (Ch. 12)   & 3--5   & Often shared with security operations. \\
Crypto (Ch. 13)          & 1--2   & Operational discipline more than ongoing engineering. \\
Orchestration (Ch. 14)   & 2--4   & Composition is the principal effort. \\
\bottomrule
\end{tabularx}
\end{table}

\paragraph{Build versus run split.} The ranges above describe
steady-state operations. The build phase --- the migration waves of
Part~III --- typically requires an additional one to two
engineer-equivalents per active wave for the duration of the wave,
on top of the steady-state figures. Institutions running multiple
waves in parallel should plan for the cumulative figure rather than
expect efficient resource sharing across waves.

%-----------------------------------------------------------------------------
\section{Scaling levers and breakpoints}
\label{sec:sz-scaling}

The reference figures imply scaling decisions at certain
breakpoints. The breakpoints below are indicative thresholds at
which the architecture would benefit from reshaping; the exact
threshold for any institution depends on its growth profile and on
its operational patterns.

\begin{itemize}[leftmargin=1.5em,nosep]
  \item \textbf{Identity above approximately fifty thousand
    subjects.} The single-cluster Keycloak deployment of the
    reference would benefit from horizontal partitioning, either
    by population (staff / student / external) or by federation
    role. Federation participation strategies may also need
    reshaping.
  \item \textbf{Audit above approximately five hundred gigabytes
    per day.} The single-cluster OpenSearch deployment of the
    reference would benefit from data-tier separation (hot indexed
    cluster, warm searchable cluster, cold object storage), and
    the analytical-store choice may need to be revisited if the
    institution had committed to a per-GB priced alternative.
  \item \textbf{Policy above approximately five thousand decisions
    per second.} The single-instance OPA deployment of the
    reference would benefit from sidecar deployment per critical
    application, with bundle distribution coordinated centrally.
  \item \textbf{Storage above the multi-petabyte threshold.} The
    single-cluster Ceph deployment of the reference would benefit
    from federation across multiple Ceph clusters with metadata
    coordination, and the operational team would need substantial
    expansion.
  \item \textbf{Multi-site institutions.} Institutions with
    multiple physical sites or with a research-cloud presence in
    multiple jurisdictions face additional design decisions
    (cross-site identity federation, cross-site policy
    consistency, cross-site audit correlation) that the
    single-site reference does not address. These would form the
    subject of a future revision of this chapter.
\end{itemize}

\noindent The breakpoints above are indicative thresholds; the
architecture remains valid beyond them, but the specific composition
of the reference instantiation may require adaptation. The
sovereignty grid of \trchap{ch:sovereignty-grid}{} continues to
apply at any scale; the per-layer instantiations do not.

\section{Down-scaled profile for smaller and teaching-led institutions}
\label{sec:sz-downscale}

The reference instantiation is sized for a research-intensive
university; it is not a floor below which the architecture ceases to
apply. The same nine layers and the same sovereignty grid hold for a
smaller or teaching-led institution, but their instantiation collapses
to a lighter profile:

\begin{itemize}[nosep]
  \item \textbf{Collapsed tiers.} The administrative, teaching and
    research planes that a large institution separates physically can
    be realised as logical segments on shared hardware; the
    distinction is enforced by policy and identity rather than by
    dedicated fabric.
  \item \textbf{Managed or mutualised compute and storage.} Where an
    institution lacks the scale to operate its own virtualisation and
    object-storage layers, these can be drawn from a national research
    cloud, a national research and education network service, or a
    managed provider under
    an appropriate data-processing agreement --- the sovereignty
    requirement constrains \emph{where} and \emph{under what terms}
    data sits, not whether the institution owns the metal.
  \item \textbf{Optional Science DMZ.} The Science~DMZ is justified by
    sustained high-volume research data transfer; a teaching-led
    institution without that traffic profile can defer it entirely
    without weakening the rest of the architecture.
  \item \textbf{Non-additive minimum team.} The role inventory of
    \S\ref{sec:sz-fte} describes responsibilities, not headcount: in a
    small institution several roles are discharged by the same person,
    and the figures there are not summed into a minimum staff count.
\end{itemize}

\noindent\textbf{Down-scaling counterpart to the breakpoints.}
Just as \S\ref{sec:sz-scaling} identifies the thresholds at which the
reference instantiation must grow, the same logic read downward
identifies what a smaller institution may legitimately omit or
mutualise: each layer degrades gracefully to a managed or shared
realisation as scale decreases, and the architecture is judged by the
sovereignty grid it satisfies, not by the size of the estate that
satisfies it. The conformance expectation for this profile is read
accordingly: the exit-confidence index of \trchap{ch:conformance}{}
(\S\ref{sec:cf-reporting}) is a maturity trajectory, and for a
mutualised layer the institution's confidence rests on the provider's
contractual exit terms rather than on a self-run test.

\paragraph{A concrete minimal-core envelope.} For planning purposes, a
teaching-led institution adopting only the minimal core --- Wave~0
governance, the Identity layer and the Observability layer, with compute
and storage mutualised as above --- can expect an indicative steady-state
envelope on the order of three to five engineer-equivalents, shared across
the roles rather than summed per layer. Mutualising compute and storage
converts their local engineering effort into provider-management overhead
rather than headcount; the Network layer collapses to logical segmentation
on shared hardware; and the optional Science~DMZ is deferred. Consistent
with the rest of this chapter, the figure is an order-of-magnitude planning
aid, not a floor or a conformance threshold: an institution that mutualises
further layers trends toward the lower bound, and the responsibilities of
\S\ref{sec:sz-fte} remain role descriptions discharged by shared staff.
